% Generated by Listas/make_statements.py from pratica-lecture-07.md, item 4. Do not edit.
\textbf{Source.} Tadelis (2013), \emph{Game Theory: An Introduction}, Exercise 12.1 ``Chicken Revisited'', p.~266 {[}PDF p.~285{]}. Reproduced verbatim. \textbf{Answer in the source:} no (none in the book; Princeton UP's website has about half, not checked).



\begin{quote}

\textbf{Chicken Revisited:} Consider the game of chicken in Section 12.2.1 with the parameters \(R = 8\), \(H = 16\), and \(L = 0\) as described there. A preacher, who knows some game theory, decides to use this model to claim that moving to a society in which all parents are lenient will have detrimental effects on the behavior of teenagers. Does equilibrium analysis support this claim? What if \(R = 8\), \(H = 0\), and \(L = 16\)?

\end{quote}



\textbf{Referred to in the statement.} The game of chicken of Section 12.2.1 (Tadelis pp.~252--254 {[}PDF p.~271--273{]}) (reproduced from the book; Figure 12.3, the game tree, is replaced by the matrix the book derives from it):



\begin{quote}

Two teenagers, players 1 and 2, have borrowed their parents' cars and decided to play the game of chicken as follows: They drive toward each other in the middle of a street, and just before impact they must simultaneously choose whether to be chicken and swerve to the right, or continue driving head on. If both are chicken then both gain no respect from their friends but suffer no losses; thus both get a payoff of 0. If \(i\) continues to drive while \(j \neq i\) plays chicken then \(i\) gains all the respect, which is a payoff of \(R\), and \(j\) gets no respect, which is worth 0. In this case both players suffer no additional losses. Finally if both continue to drive head on then they split the respect (because respect is considered to be relative), but an accident is bound to happen and they will at least be reprimanded by their parents, if not seriously injured. An accident imposes a personal loss of \(k\) on each player, so the payoff to each one is \(\frac{R}{2} - k\).



There is, however, a potential difference between these two youngsters: The punishment, \(k\), depends on the type of parents they have. For each kid, parents can be either harsh (\(H\)) or lenient (\(L\)) with equal probability, and the draws from Nature on the types of parents are independently distributed. This means that the likelihood that player \(i\)'s parents are harsh is equal to \(\frac{1}{2}\) and is independent of the type of parents that player \(j\) has. If player \(i\)'s parents are harsh then they will beat the living daylights out of their child, which imposes a high cost of an accident, denoted by \(k = H\). If instead the parents are lenient then they will give their child a long lecture on why his behavior is unacceptable, which imposes a lower cost of an accident, denoted by \(k = L < H\). Each kid knows the type of his parents but does not know the type of his opponent's parents. The distribution of types is common knowledge (this is the common prior assumption).



{[}\ldots{]} The action sets of each player are \(A_1 = \{C, D\}\) and \(A_2 = \{c, d\}\), where \(C\) (or \(c\)) stands for ``chicken'' and \(D\) (or \(d\)) stands for ``drive.''



From the extensive form in Figure 12.3 we can derive the matrix form as follows. A strategy for player 1 is denoted by \(xy \in S_1 = \{CC, CD, DC, DD\}\), where \(x\) is what he does if he is an \(L\) type and \(y\) is what he does if he is an \(H\) type. Similarly \(S_2 = \{cc, cd, dc, dd\}\). {[}\ldots{]} This results in the following matrix-form Bayesian game:



{\def\LTcaptype{none} % do not increment counter

\begin{longtable}[]{@{}lllll@{}}

\toprule\noalign{}

Player 1 ~Player 2 & \(cc\) & \(cd\) & \(dc\) & \(dd\) \\

\midrule\noalign{}

\endhead

\bottomrule\noalign{}

\endlastfoot

\(CC\) & \(0, 0\) & \(0, \frac{R}{2}\) & \(0, \frac{R}{2}\) & \(0, R\) \\

\(CD\) & \(\frac{R}{2}, 0\) & \(\frac{3R}{8} - \frac{H}{4}, \frac{3R}{8} - \frac{H}{4}\) & \(\frac{3R}{8} - \frac{H}{4}, \frac{3R}{8} - \frac{L}{4}\) & \(\frac{R}{4} - \frac{H}{2}, \frac{3R}{4} - \frac{L}{4} - \frac{H}{4}\) \\

\(DC\) & \(\frac{R}{2}, 0\) & \(\frac{3R}{8} - \frac{L}{4}, \frac{3R}{8} - \frac{H}{4}\) & \(\frac{3R}{8} - \frac{L}{4}, \frac{3R}{8} - \frac{L}{4}\) & \(\frac{R}{4} - \frac{L}{2}, \frac{3R}{4} - \frac{L}{4} - \frac{H}{4}\) \\

\(DD\) & \(R, 0\) & \(\frac{3R}{4} - \frac{L}{4} - \frac{H}{4}, \frac{R}{4} - \frac{H}{2}\) & \(\frac{3R}{4} - \frac{L}{4} - \frac{H}{4}, \frac{R}{4} - \frac{L}{2}\) & \(\frac{R}{2} - \frac{L}{2} - \frac{H}{2}, \frac{R}{2} - \frac{L}{2} - \frac{H}{2}\) \\

\end{longtable}

}



To solve for the Bayesian Nash equilibria we need to have more information about the parameters \(R\), \(H\), and \(L\). Assume that \(R = 8\), \(H = 16\), and \(L = 0\), which results in the following matrix-form game:



{\def\LTcaptype{none} % do not increment counter

\begin{longtable}[]{@{}lllll@{}}

\toprule\noalign{}

Player 1 ~Player 2 & \(cc\) & \(cd\) & \(dc\) & \(dd\) \\

\midrule\noalign{}

\endhead

\bottomrule\noalign{}

\endlastfoot

\(CC\) & \(0, 0\) & \(0, 4\) & \(0, 4\) & \(0, 8\) \\

\(CD\) & \(4, 0\) & \(-1, -1\) & \(-1, 3\) & \(-6, 2\) \\

\(DC\) & \(4, 0\) & \(3, -1\) & \(3, 3\) & \(2, 2\) \\

\(DD\) & \(8, 0\) & \(2, -6\) & \(2, 2\) & \(-4, -4\) \\

\end{longtable}

}



A quick analysis reveals that the game has a unique pure-strategy Bayesian Nash equilibrium: \((DC, dc)\).

\end{quote}

